%%=============================================================================
%% Samenvatting
%%=============================================================================

% TODO: De "abstract" of samenvatting is een kernachtige (~ 1 blz. voor een
% thesis) synthese van het document.
%
% Een goede abstract biedt een kernachtig antwoord op volgende vragen:
%
% 1. Waarover gaat de bachelorproef?
% 2. Waarom heb je er over geschreven?
% 3. Hoe heb je het onderzoek uitgevoerd?
% 4. Wat waren de resultaten? Wat blijkt uit je onderzoek?
% 5. Wat betekenen je resultaten? Wat is de relevantie voor het werkveld?
%
% Daarom bestaat een abstract uit volgende componenten:
%
% - inleiding + kaderen thema
% - probleemstelling
% - (centrale) onderzoeksvraag
% - onderzoeksdoelstelling
% - methodologie
% - resultaten (beperk tot de belangrijkste, relevant voor de onderzoeksvraag)
% - conclusies, aanbevelingen, beperkingen
%
% LET OP! Een samenvatting is GEEN voorwoord!

%%---------- Nederlandse samenvatting -----------------------------------------
%
% TODO: Als je je bachelorproef in het Engels schrijft, moet je eerst een
% Nederlandse samenvatting invoegen. Haal daarvoor onderstaande code uit
% commentaar.
% Wie zijn bachelorproef in het Nederlands schrijft, kan dit negeren, de inhoud
% wordt niet in het document ingevoegd.

\IfLanguageName{english}{%
\selectlanguage{dutch}
\chapter*{Samenvatting}
\lipsum[1-4]
\selectlanguage{english}
}{}

%%---------- Samenvatting -----------------------------------------------------
% De samenvatting in de hoofdtaal van het document

\chapter*{\IfLanguageName{dutch}{Samenvatting}{Abstract}}

Deze bachelorproef richt zich op het ontwerpen en realiseren van een datagedreven scoutingomgeving voor de BNXT League, een competitie waar clubs wel over wedstrijd- en spelersstatistieken beschikken maar waar rekruteringsbeslissingen vaak nog gebaseerd worden op intuïtie, losse bronnen en persoonlijke netwerken. De probleemstelling vertrekt vanuit deze kloof tussen beschikbare data en het ontbreken van een geïntegreerd, herhaalbaar beslissingskader voor spelersselectie, ondanks de aanwezigheid van relevante prestatie-indicatoren in bestaande basketbalstatistieken. De centrale onderzoeksvraag luidt: Hoe kan een datagedreven scoutingomgeving voor de BNXT League worden ontworpen en gerealiseerd zodat scouts spelers efficiënt kunnen analyseren en vergelijken op basis van relevante prestatie-indicatoren? 
\\\\
Het concrete doel is een proof‑of‑concept op te leveren met een centrale datalaag en een Power BI dashboard waarmee BNXT spelers gefilterd, gerangschikt en onderling vergeleken kunnen worden, aangevuld met documentatie van datamodel en indicatoren. Methodologisch wordt eerst via een literatuurstudie bepaald welke statistieken en besluitvormingskaders relevant zijn, gevolgd door het verzamelen, opschonen en modelleren van BNXT data en de implementatie van een Power BI dashboard met spelersoverzichten, detailpagina’s en vergelijkingsschermen. De omgeving wordt geëvalueerd via technische tests en een beperkte gebruiksbeoordeling aan de hand van vooraf gedefinieerde use cases. Verwacht wordt dat deze scoutingomgeving een deel van de huidige, eerder intuïtieve besluitvorming vervangt door gestructureerde en herhaalbare analyses, zodat scouts en sportief managers beter onderbouwde beslissingen kunnen nemen binnen de BNXT context.

